\section{Casimir-secular mass operator}
\label{sec:casimir}
% Status: derived/open. Eigenvalue algebra is derived; uniqueness remains an analytical subtask.

Let \(J=s(s+1)\) be the spin-Casimir eigenvalue.  The normalized two-channel operator used in this manuscript is
\begin{equation}
  \frac{Q(J)}{\mu^2}=\begin{pmatrix}
  0 & \sqrt{J}\\
  \sqrt{J} & -J
  \end{pmatrix}.
  \label{eq:QJ}
\end{equation}
Its characteristic equation is
\begin{equation}
  \det\left(xI-\frac{Q(J)}{\mu^2}\right)=x^2+Jx-J=0.
  \label{eq:charpoly}
\end{equation}
The two roots are
\begin{equation}
  x_\pm(J)=\frac{-J\pm\sqrt{J^2+4J}}{2}.
  \label{eq:roots}
\end{equation}
They obey
\begin{equation}
  x_+(J)+x_-(J)=-J,
  \qquad
  x_+(J)x_-(J)=-J.
  \label{eq:root-identities}
\end{equation}
For \(J>0\), the positive branch is positive and the negative branch is negative.

Historically, the electroweak number is obtained by sampling the positive branch at the two Casimir values
\begin{equation}
  J_H=\frac34,
  \qquad
  J_{\rm adj}=2.
  \label{eq:historical-J-samples}
\end{equation}
This gives
\begin{equation}
  \sdV = 1-\frac{x_+(3/4)}{x_+(2)}.
  \label{eq:sdvdef}
\end{equation}
The expression is dimensionless; the overall scale \(\mu\) cancels.  The cancellation is essential for an electroweak interpretation because Eq.~\eqref{eq:sdvdef} can then be compared with a projective W/Z mass ratio.  In the electroweak reading developed here, \(J_H=3/4\) is the Higgs/order-parameter doublet invariant and \(J_{\rm adj}=2\) is the adjoint/current invariant.  The ordered use of these two samples in the W/Z quotient remains the central derivation problem.

The relation to the original Rivero operator is obtained by restoring the Poincare Casimir eigenvalues \(C_1=m^2\) and \(C_2=-m^2s(s+1)\).  In the original construction, a mass-squared operator is formed from combinations of Casimir invariants, with a high-spin condition designed to preserve the asymptotic mass behavior relevant for Regge trajectories \cite{Rivero2006}.  Equation~\eqref{eq:QJ} is the reduced dimensionless representative of that branch problem.

\subsection{What the algebra fixes}

The operator in Eq.~\eqref{eq:QJ} fixes three algebraic facts.  First, the trace is \(-J\).  Second, the determinant is \(-J\).  Third, the two eigenvalues are tied by the same invariant:
\begin{equation}
  \operatorname{tr}\frac{Q(J)}{\mu^2}=-J,\qquad
  \det\frac{Q(J)}{\mu^2}=-J.
  \label{eq:trace-det-fixed}
\end{equation}
Thus a physical derivation has to explain a simultaneous trace and determinant assignment.  A model that gives only a square-root branch, or only a \(J\)-dependent diagonal term, leaves the DeVries construction incomplete.  The required two-channel matrix has the schematic form
\begin{equation}
  \mathcal M_J=
  \begin{pmatrix}
  0 & \kappa_J\\
  \kappa_J & -\tau_J
  \end{pmatrix},
  \qquad
  \kappa_J^2=J,\qquad \tau_J=J.
  \label{eq:minimal-two-channel}
\end{equation}
This is the compact target for the endpoint, Kaluza--Klein, and \(G_2\) sections.

The square-root radical is then the discriminant of this two-channel problem:
\begin{equation}
  \Delta_J=J^2+4J.
  \label{eq:discriminant}
\end{equation}
The radical has to attach to a physical object.  The radical-placement section lists the main options.  This paper follows the pole-spectrum placement because the positive and negative branches are interpreted as a single low-energy spectral clue.

\subsection{Minimal assumptions and theorem target}

A complete derivation should state the minimal assumptions that select Eq.~\eqref{eq:QJ}.  The current minimal list is:
\begin{enumerate}
\item dependence on available Poincare or internal invariants;
\item dimensional consistency after restoring the mass scale \(\mu\);
\item a two-branch secular structure;
\item preservation of the high-spin or Regge asymptotic datum in the originating problem;
\item a trace/determinant condition equivalent to Eq.~\eqref{eq:trace-det-fixed}.
\end{enumerate}
These assumptions define the uniqueness target.  The theorem to prove would read: within a stated class of two-channel operators built from the allowed invariants, the high-spin condition and trace/determinant constraints select Eq.~\eqref{eq:QJ} up to similarity transformations and normalization.

This theorem would also clarify the status of the negative branch.  Since \(x_-(J)\) is fixed by the same trace and determinant as \(x_+(J)\), a derivation of Eq.~\eqref{eq:QJ} automatically derives both branches.  The electroweak interpretation can then ask which fields or boundary data realize the two eigenvectors.
